\documentclass{tongji-comporg-report}

\ReportTopic{乘法器}
\ReportStudentID{}
\ReportStudentName{}
\ReportMajor{}
\ReportTeacher{}

\begin{document}

\MakeReportCover

\begin{ReportBody}

\ReportSection{实验环境与实验内容}
实验内容：使用硬件描述语言verilog，设计实现32 位无符号乘法器和32 位有符号乘法器。

\ReportSectionGap

\ReportSection{乘法器的结构框架与分析解释}

\ReportSectionGap

\ReportSection{实验仿真过程及波形分析}

\ReportSectionGap

\ReportSection{实验验证}

\ReportBigGap[1.65cm]

\ReportSection{总结与体会}

\ReportBigGap[3.55cm]

\ReportSection{附件（所有源码程序）}

\end{ReportBody}

\end{document}
